We study numerical adaptations and low-level implementations of a fixed random feature Galerkin solver for incompressible flow, separating the effects of the trial space, time schedule, operator cost, and runtime overhead. For two-dimensional flow in a finite annulus without imposed reflection symmetry, we specify boundary-compatible features, weighted rank truncation, and a weak discretization with an independent mean-flow block. We derive a variable-step SBDF2 bridge with boundary lifting and establish its local consistency. Common coordinates for the time-phase matrices of each angular mode yield a geometric anchor that minimizes a stated worst-case condition-number bound. The implementation analysis describes FP64 FFT packing and shared-memory fusion, warp-level probe reduction, native C radial kernels, and setup and compilation-artifact reuse. At $Re=100$ and $T=72$, all 78 GPU trajectories, comprising 39 configurations under two cache policies, meet the prescribed $0.5\%$ velocity and vorticity tolerances. Piecewise schedules remove $11.11\%$ and $9.36\%$ of true nonlinear evaluations. FFT fusion reduces short-trajectory marching time by $9\%$--$17\%$ but increases the first full-solve cost by $23\%$--$69\%$. Twenty new CPU trajectories supplement feature and switch controls. Both fixed-midpoint frequency configurations qualify with analytic layers, while all 12 configurations without nonzero-mode layers fail; retained rank also decreases, so this is not an equal-rank enrichment ablation. Delaying the switch only modestly improves accuracy at greater nonlinear cost, and radial/modal profiles distinguish the distribution of velocity and vorticity errors. Reduced Legendre spaces also qualify, and the final isolated timing campaign is not accepted; an overall speed advantage over comparably optimized traditional CFD therefore remains unestablished. The contribution is a reproducible combination of discrete derivations, source-level cost explanations, and scoped positive and negative evidence. Local consistency, transformed conditioning, and numerical qualification are not interpreted as guarantees of nonlinear stability or universal performance.
